La respuesta inicial de la IA fue sintácticamente correcta en Mermaid y recogió adecuadamente la estructura general del flujo (solicitud de aplicación de cupón, consulta en base de datos y bloque \texttt{alt/else/end}). Sin embargo, fue necesario realizar tres correcciones técnicas para alinear el diagrama con el resto de la arquitectura:
\begin{itemize}[leftmargin=*,nosep]
\item \textbf{Método HTTP y ruta REST distintos al glosario común:} La IA propuso \texttt{POST /api/cart/apply-coupon (codigo, idCarrito)}. Se corrigió para emplear el método y ruta oficiales del glosario común: \texttt{PUT /api/carrito/cupon (codigo)}.
\item \textbf{Incoherencia con las entidades \texttt{CARRITO} y \texttt{CUPON} del diagrama ER:} La IA ejecutó \texttt{UPDATE carts SET coupon\_id = idCupon, discount\_amount = descuento}, cuando en \texttt{er.mmd} la clave primaria de \texttt{CUPON} es \texttt{codigo} y \texttt{CARRITO} solo referencia \texttt{cupon\_codigo FK}.
\item \textbf{Correspondencia con los métodos de las clases \texttt{Cupon} y \texttt{Carrito}:} Se reemplazaron las descripciones textuales del Backend por los métodos explícitos de \texttt{clases.mmd}: \texttt{findCuponByCodigo(codigo)}, \texttt{Cupon.esValido()}, \texttt{Carrito.aplicarCupon(codigo)}, \texttt{Cupon.calcularDescuento(subtotal)} y \texttt{Carrito.calcularTotal()}, precisando el estado \texttt{422 Unprocessable Entity} para cupones inválidos e incluyendo \texttt{mirrorActors: false}.
\end{itemize}
